% ============================================
% Johns Hopkins Letter Template
% Assistant Professor Cover Letter – Denison University
% ============================================

\documentclass[11pt,letterpaper]{letter}

\usepackage{graphicx}
\usepackage{eso-pic}
\usepackage{geometry}
\usepackage{microtype}
\usepackage[hidelinks]{hyperref}

% ----- Page Setup -----
\geometry{left=25mm,right=25mm,top=25mm,bottom=25mm}

% ----- Logo Placement (first page only) -----
\newcommand{\PutJHULogoFG}[1]{%
  \AddToShipoutPictureFG*{%
    \AtPageUpperLeft{%
      \hspace{10mm}%
      \raisebox{-38mm}{%
        \includegraphics[width=6cm]{#1}%
      }%
    }%
  }%
}

% ----- Sender Information -----
\signature{Decory Edwards}
\address{Department of Economics \\ Johns Hopkins University \\ Baltimore, MD 21218 \\ \texttt{dedwar65@jh.edu}}

\begin{document}
\PutJHULogoFG{Images/jhulogo.png}

\begin{letter}{Search Committee \\
Department of Economics \\
Denison University \\
Granville, OH 43023 \\ USA}

\opening{Dear Members of the Search Committee,}

I am writing to express my interest in the tenure-track Assistant Professor position in Finance in the Department of Economics at Denison University, beginning August 2026. I am a Ph.D. candidate in Economics at Johns Hopkins University, advised by Professors Christopher D.~Carroll, Nicholas W.~Papageorge, and Stelios Fourakis, and I expect to receive my degree in May 2026. My research fields are macroeconomics, wealth and income dynamics, and computational economics.

My research agenda focuses on understanding the determinants of wealth inequality and the mechanisms through which heterogeneity in returns to wealth shapes aggregate outcomes. In my job-market paper, I estimate the degree of heterogeneity in individual portfolio returns using micro-data from the Survey of Consumer Finances and simulate a structural model in which households face idiosyncratic return risk. I show that allowing for persistent heterogeneity in returns substantially improves the model’s ability to match the empirical distribution of wealth, offering a tractable way to connect micro-level return dispersion to macroeconomic inequality.

In a second project, I use the Health and Retirement Study to examine the relationship between interpersonal trust and portfolio performance. Building on the literature that finds a hump-shaped relationship between trust and income, I show that a similar relationship holds for individual returns to wealth measured in the HRS data, highlighting a behavioral channel through which social attitudes shape saving and investment outcomes. This work also provides new estimates of the persistent component of returns to net worth, which motivate the calibration of heterogeneity in my structural model.

Denison’s commitment to liberal arts education, interdisciplinary exploration, and open inquiry aligns strongly with my approach to research and teaching. I am excited about the department’s expansion to include a new B.S. in Finance and about contributing to a program that values context, ethical reasoning, and critical thinking in the study of financial markets. The opportunity to teach \textit{Principles of Finance}, the Finance capstone, introductory and intermediate economic theory, and finance electives fits naturally with my background in macroeconomics, portfolio behavior, and quantitative modeling. I am particularly drawn to Denison’s emphasis on active learning, debate, and constructive engagement across different perspectives—values that strongly shape my own teaching philosophy.

\textbf{Teaching.} My teaching experience centers on undergraduate macroeconomics and an upper-level macro policy seminar, where I have led weekly discussion sections and held regular office hours for classes ranging from 16 to 40 students. My approach is student-centered: I aim to help students “convince themselves” that they have moved from confusion to understanding by holding structured office-hour coaching that builds trust and confidence. At Denison, I am prepared to teach \textit{Principles of Finance}, the Finance \textit{Capstone Seminar}, \textit{Principles of Macroeconomics}, \textit{Intermediate Macroeconomic Theory}, and finance electives aligned with students’ interests. I also welcome the opportunity to mentor undergraduate research, support experiential learning, and contribute to advising across the major and minor.

I believe I would be a strong fit because my empirical and computational experience allows me to (i) teach finance and economics using real-world datasets and reproducible methods, (ii) create active learning environments that emphasize analytical reasoning and ethical decision-making, and (iii) contribute meaningfully to a collaborative department that values diverse scholarly approaches. I am excited about Denison’s investment in faculty research, including pre-tenure leave and opportunities for student-faculty collaboration, and I look forward to contributing to the department’s continued growth and excellence.

I would be honored to join Denison University as part of a community of teacher-scholars dedicated to undergraduate education, vibrant debate, and high-quality research. Thank you for your consideration; I welcome the opportunity to discuss my fit with your department at your convenience.

\closing{Sincerely,}

\end{letter}
\end{document}
